% sections/7_technical.tex
\chapter{Техническое описание прохождения маршрута}
\label{ch:technical}

В этой главе даётся подробное описание каждого ходового дня: характер пути, ориентиры, действия группы, погодные условия, места ночёвок и другие детали.

% ----------------------------------------------------------------------
% День 1
% ----------------------------------------------------------------------
\dayentry[day1]{День 1. 16.07.2022. Пос. Турасу — дол. р. Тон}

\begin{tourinfotable}
\tourinfo{Дата}{16.07.2022}
\tourinfo{День пути}{1}
\tourinfo{Участок}{Пос. Турасу — дол. р. Тон}
\tourinfo{Протяжённость (км)}{8,5}
\tourinfo{ЧХВ}{3:15}
\tourinfo{Набор / сброс высоты}{+450 / –20 м}
\tourinfo{Максимальная высота (м)}{1800}
\tourinfo{Погода}{Ясно, +25°C, ветер слабый}
\tourinfo{Характер пути}{Лесная тропа, подъём по серпантину}
\tourinfo{Ночёвка}{Поляна на левом берегу р. Тон, есть вода}
\end{tourinfotable}

\textbf{Подробное описание:}
Выехали из посёлка Турасу в 9:00. Тропа хорошо набита, идёт по лесу с постепенным набором высоты. Пересекли несколько ручьёв. В 12:15 остановились на обед. Во второй половине дня вышли к реке Тон, нашли удобное место для лагеря. Палатки поставили на поляне, вода рядом.

% ----------------------------------------------------------------------
% День 2
% ----------------------------------------------------------------------
\dayentry[day2]{День 2. 17.07.2022. дол. р. Тон — дол. р. Зындан}

\begin{tourinfotable}
\tourinfo{Дата}{17.07.2022}
\tourinfo{День пути}{2}
\tourinfo{Участок}{дол. р. Тон — дол. р. Зындан}
\tourinfo{Протяжённость (км)}{12,2}
\tourinfo{ЧХВ}{5:00}
\tourinfo{Набор / сброс высоты}{+620 / –150 м}
\tourinfo{Максимальная высота (м)}{2100}
\tourinfo{Погода}{Переменная облачность, +20°C}
\tourinfo{Характер пути}{Травянистые склоны, местами осыпи}
\tourinfo{Ночёвка}{Сухая терраса у слияния ручьёв}
\end{tourinfotable}

\textbf{Подробное описание:}
Вышли в 8:30. Тропа местами теряется, ориентировались по каменным турам. На перевальной седловине открывается вид на долину Зындан. Спуск по осыпи средней крупности. В 15:00 вышли к месту ночёвки. Вода из ручья.

% ----------------------------------------------------------------------
% Добавляйте следующие дни по аналогии, копируя блок с \dayentry и таблицей
% ----------------------------------------------------------------------

\dayentry[day3]{День 3. 18.07.2022. дол. р. Зындан — пер. 267}

\begin{tourinfotable}
\tourinfo{Дата}{18.07.2022}
\tourinfo{День пути}{3}
\tourinfo{Участок}{дол. р. Зындан — пер. 267}
\tourinfo{Протяжённость (км)}{9,0}
\tourinfo{ЧХВ}{6:30}
\tourinfo{Набор / сброс высоты}{+870 / –50 м}
\tourinfo{Максимальная высота (м)}{4086}
\tourinfo{Погода}{Облачно, ветер, временами снег}
\tourinfo{Характер пути}{Крупная осыпь 20–25°, снежник}
\tourinfo{Ночёвка}{На перевале (под седловиной)}
\end{tourinfotable}

\textbf{Подробное описание:}
Подъём занял большую часть дня. Осыпь подвижная, двигались связками. На снежнике организовали перила для страховки. Верхняя точка — седловина перевала 267 (1Б, 4086 м). Из-за ухудшения погоды решили не спускаться, а заночевать на полке под перевалом.

% Если необходимо вставить фото или схему дня, используйте команду \insertmap
% \insertmap[Фото перевала 267]{figures/pass267.jpg}